% page 461 of the facsimile (image p-460 of the scan)
% block 170.2 x 237.7 mm ; left margin 31.8 mm
\pgc{10.160mm}{0.000mm}{
\l{\cel{51}453}
\l{}
\l{\sou{populo}. Ensemblo ek la personi qui habitas la sama lando, esas}
\l{\cel{3}(maxim-multa-kaze) sama-rasa, ed uzas la sama linguo. - II.}
\l{\cel{3}La korpo di la naciono, opoze a la suvereno ed a la nobeli.}
\l{\cel{3}- EFIS.}
\l{}
\l{\sou{populara}. I. Qua esas komprenebla (kom texto) da la personi or-}
\l{\cel{3}dinara, ne-specalista, ne-ciencista. - II. Pri qua la populo}
\l{\cel{3}opinionas favoroze. - DEFIRS.}
\l{}
\l{\sou{popurito}. Verko muzikala, literaturala, ek texti de origini di-}
\l{\cel{3}versa. - DEF.}
\l{}
\l{\sou{por}.\cel{2}Profite da... Destine di...}
\l{}
\l{\sou{poro}. I. Singla ek la orifici preske neperceptebla di la pelo di}
\l{\cel{3}animalo, per qui eventas la transpiro. - II. Singla ek la}
\l{\cel{3}orifici analoga,ye qui truizesas la vejetanti. - III. Singla}
\l{\cel{3}ek la interstici qui interseparas la molekuli di korpo ed igas}
\l{\cel{3}ica plu o min permeebla.- DEFIRS.}
\l{}
\l{\sou{porcelano}. Ceramikajo blanka, tre dina, importacita de la esto-}
\l{\cel{3}regioni di Azia (precipue Chinia) aden Europa en la yarcento}
\l{\cel{3}16-esma. - DEFIS.}
\l{}
\l{\sou{porciono}. Parto qua destinesas a singlu, lor disdono o distributo.}
\l{\cel{3}DEFIRS.}
\l{}
\l{\sou{pordo}. I. Aperturo facita en la muro di klozajo, di edifico, di}
\l{\cel{3}chambro, por povar enirar od ekirar. - II. Klapo qua klozas}
\l{\cel{3}ta aperturo e povas apertesar por paso. - EFIS.}
\l{}
\l{\sou{porelo}. (\sou{bot.}) Supo-planto, di la genero \textquotedbl{}alio\textquotedbl{}. - L.\cel{1}\sou{allium}}
\l{\cel{3}\sou{porrum}. - DFIRS.}
\l{}
\l{\sou{porfiro}. (\sou{mineral}.) Roko harda, obskure-reda, en qua dissemesas}
\l{\cel{3}makuli blanka, quan la antiqui obtenis de la regioni alte-situit}
\l{\cel{3}di Egiptia. - DEFIS.}
\l{}
\l{\sou{porismo}. (\sou{geom.)}\cel{2}(\sou{anti}que).Teoremo nekompleta, qua establisas}
\l{\cel{3}proprajo komuna inter ula elementi quin la enunco determinas,}
\l{\cel{3}ed altra teoremi qui rezultas de la hipotezo. - DEFI.}
\l{}
\l{\sou{porko}. (\sou{zool.}) Mamifero domestika quan onu repletigas por manjor}
\l{\cel{3}lu. - L. sus (scrofa domestica).- EFIS.}
\l{}
\l{\sou{porkespino}. (\sou{zool.}) Mamifero rodera di qua la korpo havas kom}
\l{\cel{3}armi pikili. - L.\cel{1}\sou{hystrix}.}
\l{}
\l{\sou{pornografo}. I. Autoro qua kompozas texti pri prostituco. - II.}
\l{\cel{3}Autoro, piktisto, obcena. - DEFIS.}
\l{}
\l{\sou{pornografio}. I. Traktato pri prostituco. - II. Verki literatura,}
\l{\cel{3}arta, quin - kontre-etike - onu igas servar libertinismo.}
}
